\chapter{Bursal Aspiration and Injection}

\section*{Overview}
Bursal aspiration and injection is a procedure in which fluid is removed from an inflamed bursa and, when appropriate, a corticosteroid is injected into the bursa to relieve inflammation and pain. It is most commonly performed for bursitis of the shoulder, elbow, hip, or knee.
\section*{Alternative Names}
Bursa aspiration; bursa injection; aspiration and injection of a bursa; corticosteroid bursal injection
\section*{Principle}
Aspiration removes excess fluid and provides material for analysis or symptom relief, while corticosteroid injection reduces local inflammation of the bursal lining, decreasing swelling and pain.
\section*{History}
Aspiration and injection of inflamed bursae has been practiced since the early 20th century, and became routine with the availability of injectable corticosteroids after the 1950s.
\section*{Why It Is Done}
Ordered for painful bursitis with significant effusion, to relieve symptoms, to obtain fluid for testing when septic bursitis or gout is suspected, and to deliver corticosteroid directly to the inflamed bursa.
\section*{Is This a Primary or Secondary Test or Procedure}
A primary therapeutic and diagnostic procedure for bursitis.
\section*{How It Is Performed}
The bursa is palpated and the skin prepared aseptically. A needle is inserted into the bursal cavity, fluid is aspirated, and corticosteroid may then be injected through the same needle. Ultrasound guidance may be used for deeper bursae.
\section*{Preparation}
Informed consent and review of the patient's anticoagulation status. The injection is deferred if there is overlying cellulitis or suspected infection.
\section*{Contraindications and Precautions}
Contraindications include infection in the area and bleeding diathesis. Precautions include repeated corticosteroid injections, which can weaken tendons, and use near the Achilles or patellar tendon.
\section*{Risks}
Pain, transient flare of symptoms, infection, bleeding, skin atrophy at the injection site, and tendon rupture after repeated injections.
\section*{Aftercare and Recovery}
The site is bandaged and the patient rests the area for a day. Pain relief from the corticosteroid typically begins within days, while fluid aspirated for testing is sent to the laboratory.
\section*{Outcome and Follow-up}
Improvement in pain and swelling indicates a good response. Fluid analysis results determine whether the bursitis is inflammatory, crystalline (gout or pseudogout), or infectious.
\section*{Expected Results}
Relief of symptoms and a benign fluid profile with no crystals or organisms.
\section*{Follow-up}
Patients are reassessed within weeks; persistent symptoms may prompt imaging, repeat aspiration, or further treatment including physical therapy.
\section*{References}
\begin{enumerate}
  \item MedlinePlus. Bursitis. U.S. National Library of Medicine; 2025.
  \item Current Medical Diagnosis and Treatment. McGraw-Hill; 2025.
\end{enumerate}

\clearpage
